\section*{Введение}
\addcontentsline{toc}{section}{Введение}

В главе~1 к полю присоединялся \emph{один} корень неприводимого многочлена (поле корня). Здесь мы присоединяем
\emph{все} корни многочлена и изучаем гомоморфизмы получающихся полей~--- это подготовка к теории Галуа (глава~3), где
главным объектом станет группа автоморфизмов поля разложения. Основной инструмент по-прежнему элементарен: деление
многочленов с остатком, однозначность разложения в $F[X]$ и мультипликативность степеней.

\begin{itemize}[leftmargin=*]
\item \emph{Гомоморфизмы простых расширений} (\rf{2.1}--\rf{2.4d}). $F$-гомоморфизмы $F[\alpha]\to\Om$ отвечают корням
минимального многочлена $\alpha$ в $\Om$ (без учёта кратностей), а продолжения гомоморфизма $\varphi_0\colon F\to\Om$~---
корням многочлена $\varphi_0f$. Неприводимость $X^p-a$ при $a\notin F^p$ и единственность гомоморфизма из его поля корня;
неприводимость $X^{p^n}-a$ (\rf{2.3a}).
Присоединение второго корня: результат зависит от того, корнем какого неприводимого множителя над $F[\alpha]$ является
второй корень (примеры $X^4-2$, $X^3-2$, многочлен степени $9$ с тремя неизоморфными полями, многочлен степени $10$ с
корнями $x_i+x_j$, где $x_i$~--- корни $X^5-X-1$, и многочлен $X^5-5X+12$ с группой $D_5$, у которого разные множители
дают одно и то же поле).
\item \emph{Поля разложения} (\rf{2.4e}--\rf{2.11}). Существование поля разложения и оценка $[E_f:F]\le(\deg f)!$; поля
разложения многочленов второй и третьей степени, кругового многочлена $\Phi_p$, многочлена Артина--Шрайера $X^p-X-a$ и
двучлена $X^n-a$. Подробно разобрано, когда $F[\sqrt[n]a]$~--- поле разложения $X^n-a$ (контрпримеры $X^4+1$ над $\Q$ и
$X^9-2$ над $\F_7$, случаи, когда это влечёт $\mu_n\subset F$, и роль неприводимости: $X^5-32$); попутно: $X^4+1$
неприводим над $\Q$, но приводим по модулю каждого простого. Найдено поле разложения $X^m-1$ над $\F_p$: его
степень равна порядку $p$ по модулю наибольшего делителя $m$, взаимно простого с $p$.
\item \emph{Гомоморфизмы алгебраических расширений} (\rf{2.12}--\rf{2.16}). Главная оценка этой главы: число $F$-гомоморфизмов
расширения, порождённого корнями $f$, в поле, расщепляющее $f$, не больше степени и равно ей при отсутствии кратных
корней; единственность поля разложения с точностью до изоморфизма (и отсутствие выделенного изоморфизма); оценка
$|\Hom_F(E,L)|\le[E:F]$ для любого $L$; общее конечное расширение, содержащее копии нескольких данных полей, и когда
такое общее расширение универсально (кодекартов квадрат, \rf{2.15b}).
\item \emph{Кратные корни, сепарабельность, совершенные поля} (\rf{2.17}--\rf{2.26}). Наибольший общий делитель и
делимость многочленов не меняются при расширении поля; кратность корня не зависит от выбора поля; кратный корень~---
общий корень $f$ и $f'$; неприводимый многочлен имеет кратные корни лишь в характеристике $p$ и лишь если он многочлен
от $X^p$, и тогда все его корни кратные. Сепарабельные многочлены (критерий $\gcd(f,f')=1$) и неприводимость многочлена
Артина--Шрайера $X^p-X-a$ (\rf{2.22b}); попутно~--- формулы Виета для $X^p-X$ и два доказательства
теоремы Вильсона (\rf{2.22c}). Сепарабельные элементы и расширения: критерий $E=F[E^p]$, сепарабельность
расширения, порождённого сепарабельными элементами (два доказательства: через критерий и подсчётом гомоморфизмов), и
транзитивность сепарабельности для конечных расширений (\rf{2.22d}--\rf{2.22h}; общий случай~--- в главе~6). Совершенные поля: критерий <<поле совершенно тогда и только тогда,
когда каждый неприводимый многочлен сепарабелен>>, или когда каждое его алгебраическое расширение сепарабельно
(\rf{2.24a}); примеры (характеристика $0$, конечные поля, поля, алгебраические над $\F_p$, алгебраически замкнутые поля;
несовершенные поля: $\F_p(T)$, поля рациональных функций и формальных рядов Лорана над полем характеристики $p$); алгебраическое расширение совершенного
поля совершенно (\rf{2.25a}); совершенное поле $\F_p(T^{1/p^\infty})$~--- объединение возрастающей цепочки несовершенных
(\rf{2.25b}). В заключение~--- сводка о связи (не)различия корней и (не)приводимости (\rf{2.26a}).
\item \emph{Упражнения} 2-1--2-6 с решениями: квадратичные расширения и подгруппы $F^\times/F^{\times2}$ (попарная
неизоморфность $\Q[\sqrt p]$, единственность поля из $p^2$ элементов); приводимый многочлен $X^p-X-a$ раскладывается
на линейные множители, а $X^p-X-1$ неприводим над $\Q$; поле разложения $X^5-2$ над $\Q$ (степень $20$); поля
разложения $X^{p^m}-1$ и $X^{p^m}-X$ над $\F_p$ (единственность поля из $p^m$ элементов); многочлен $f/\gcd(f,f')$;
запись $f=g(X^{p^e})$ неприводимого многочлена с неприводимым сепарабельным $g$.
\end{itemize}

Всюду в этой главе $F$~--- поле; <<расширение $F$>>~--- поле, содержащее $F$; $F$-гомоморфизм полей тождествен на $F$.
